\documentclass[11pt,a4paper]{article}

% --- Packages ---
\usepackage[utf8]{inputenc}
\usepackage[english]{babel}
\usepackage{amsmath, amssymb, amsfonts, bm}
\usepackage{geometry}
\usepackage{graphicx}
\usepackage{subcaption}
\usepackage{booktabs}
\usepackage{hyperref}
\usepackage{cite}
\usepackage{float}
\usepackage{xcolor}
\usepackage{enumitem}

\geometry{a4paper, top=2.5cm, bottom=2.5cm, left=2.5cm, right=2.5cm}

\hypersetup{
    colorlinks=true,
    linkcolor=blue!70!black,
    citecolor=green!50!black,
    urlcolor=cyan!70!black
}

% --- Macros ---
\newcommand{\bq}{\bm{q}}
\newcommand{\bdq}{\dot{\bm{q}}}
\newcommand{\bddq}{\ddot{\bm{q}}}
\newcommand{\btau}{\bm{\tau}}
\newcommand{\bM}{\bm{M}}
\newcommand{\bC}{\bm{C}}
\newcommand{\bG}{\bm{G}}
\newcommand{\bB}{\bm{B}}
\newcommand{\Real}{\mathbb{R}}
\newcommand{\dgr}{^{\circ}}

% ============================================================
\begin{document}

\begin{titlepage}
    \centering
    \vspace*{2cm}
    {\Large\bfseries Field and Service Robotics --- Final Technical Report\par}
    \vspace{1.5cm}
    {\LARGE\bfseries Modeling and Multi-Strategy Control of a\\[6pt]
    Variable-Length Wheeled Inverted Pendulum Robot\par}
    \vspace{1cm}
    {\large RoTino: a Wheeled Biped Balancing Platform\par}
    \vspace{2cm}
    {\large Aldo Rosicarelli\par}
    \vspace{0.5cm}
    {\large Academic Year 2025--2026\par}
    \vfill
    {\large Field and Service Robotics\par}
\end{titlepage}

% ============================================================
\tableofcontents
\newpage

% ============================================================
\section{Introduction}
\label{sec:intro}

Wheeled biped robots occupy a particularly interesting niche in the landscape of field and
service robotics: they combine the high mobility and energy efficiency of wheeled platforms
with the postural versatility of legged systems.  By adjusting their centre of mass (CoM)
height through variable leg length, they can negotiate uneven terrain, step over obstacles,
or adopt a compact stance for confined environments.  The archetype is the family of
\emph{Variable-Length Wheeled Inverted Pendulum} (VL-WIP) robots --- two-wheeled
self-balancing machines whose legs can change length to modulate the height at which the
upper-body mass is carried.

The control problem is inherently challenging because the system is \emph{underactuated}:
the pitch angle of the torso is not directly actuated, but must be controlled indirectly
through the centroidal coupling with the wheel acceleration.  This is analogous to the
classical inverted-pendulum-on-a-cart problem, but with the additional complexity of a
multi-link leg mechanism, configuration-dependent inertia, and the variable pendulum length
that makes the dynamics time-varying even at equilibrium.

The project selected for this course is \textbf{RoTino}, a simulated wheeled biped modelled
after the platform described by Cui et al.~\cite{cui2022}.  The development follows a
two-stage methodology:

\begin{enumerate}[leftmargin=*]
    \item \textbf{MATLAB/Simulink prototyping.}  The Euler--Lagrange dynamics are derived
          symbolically, and three control laws --- cascade PID, LQR + MPC, and super-twisting
          SMC --- are implemented and validated in ODE45/Simulink simulations of the planar
          model.
    \item \textbf{ROS~2 / Gazebo validation.}  The controllers are ported to Python ROS~2
          nodes and tested in closed-loop Gazebo simulation of a fully articulated 3D URDF
          robot model, including sensor noise, actuator dynamics, and contact physics.
\end{enumerate}

The choice is motivated by several factors:
\begin{itemize}[leftmargin=*]
    \item \textbf{Rich nonlinear underactuated dynamics.}  A four-DOF planar chain with only
          three actuators requires careful treatment of the actuation matrix and its
          singularities.
    \item \textbf{Cascade control hierarchy.}  The natural time-scale separation between
          balance ($\sim$500~Hz) and locomotion ($\sim$100~Hz) makes the system an ideal
          testbed for hierarchical control.
    \item \textbf{Comparison across control paradigms.}  Three qualitatively different
          controllers are benchmarked on the same model, enabling a principled comparison of
          performance, robustness, and implementation complexity.
    \item \textbf{Software-in-the-loop validation.}  The full pipeline from symbolic
          derivation to closed-loop ROS~2 simulation bridges theory and practice.
\end{itemize}

The report is structured as follows.  Section~\ref{sec:topology} analyses the topology and
degrees of freedom.  Section~\ref{sec:model} derives the dynamic model.
Sections~\ref{sec:pid}, \ref{sec:lqr_mpc}, and~\ref{sec:smc} describe the three control
strategies.  Section~\ref{sec:framework} presents the simulation and benchmarking framework.
Section~\ref{sec:results} discusses the results.  Section~\ref{sec:conclusions} concludes.


% ============================================================
\section{Topology Analysis and Degrees of Freedom}
\label{sec:topology}

\subsection{Mechanical Architecture}

The robot (Figure~\ref{fig:architecture}) is a planar kinematic chain comprising four rigid
bodies:

\begin{enumerate}
    \item \textbf{Wheel} ($r = 0.127$~m, $m_w = 3.5$~kg, $I_w = 0.1$~kg\,m$^2$): rolling
          without slipping on the ground.
    \item \textbf{Shank link} ($l_1 = 0.45$~m, $m_1 = 1.2$~kg): connects the wheel axle to
          the knee joint.
    \item \textbf{Thigh link} ($l_2 = 0.45$~m, $m_2 = 5.3$~kg): connects the knee to the
          hip joint.
    \item \textbf{Torso} ($l_3 = 0.35$~m, $m_3 = 60.0$~kg, wheel track $d = 0.63$~m): the
          upper body, connected to the hip.
\end{enumerate}

All link inertias are modelled as slender rods: $I_k = \frac{1}{12}m_k l_k^2$.

\begin{figure}[H]
    \centering
    \includegraphics[width=0.38\textwidth]{figures/mechanical_architecture.png}
    \caption{Schematic of the RoTino platform: torso, hip joint, ballast, thigh link, knee
    joint, shank link, wheel and wheel joint.}
    \label{fig:architecture}
\end{figure}

\subsection{Gr\"ubler Constraint and DOF Count}

A planar system of $n_b = 4$ bodies has $3 n_b = 12$ unconstrained DOFs.
Table~\ref{tab:dof} summarises the constraints that reduce this count to~4.

\begin{table}[H]
    \centering
    \caption{Constraint analysis for the planar RoTino chain.}
    \label{tab:dof}
    \begin{tabular}{llcc}
        \toprule
        Constraint & Description & Equations & Holonomic? \\
        \midrule
        Wheel--ground contact & Fixed contact height $y = r$ & 1 & Yes \\
        Pure rolling & $\dot{x} = r\,\dot{\phi}_w$ & 1 & No (Pfaffian) \\
        Wheel--shank revolute & 2 translational DOFs locked & 2 & Yes \\
        Shank--thigh revolute (knee) & 2 translational DOFs locked & 2 & Yes \\
        Thigh--torso revolute (hip) & 2 translational DOFs locked & 2 & Yes \\
        \midrule
        \textbf{Total constraints} & & \textbf{8} & \\
        \textbf{Resulting DOF} & $12 - 8$ & \textbf{4} & \\
        \bottomrule
    \end{tabular}
\end{table}

The generalised coordinates, following Cui et al.\ (Table~1), are:
\begin{equation}
    \bq = \begin{bmatrix} q_w & q_1 & q_2 & q_3 \end{bmatrix}^T
\end{equation}
where $q_w$ is the wheel rotation angle (encoding horizontal displacement $s = R_w q_w$
via rolling), $q_1$ is the absolute shank angle from the vertical, $q_2$ the absolute thigh
angle, and $q_3$ the absolute torso angle.

\subsection{Full 3D Model: Seven-DOF Formulation and Null-Space Projection}

For the 3D Gazebo simulation, the full robot is described by $n = 7$ generalised
coordinates:
\begin{equation}
    \bq_{3D} = [x,\ y,\ \psi,\ \theta,\ l,\ \phi_R,\ \phi_L]^T
\end{equation}
where $(x,y)$ is the axle midpoint on the ground plane, $\psi$ is the heading (yaw),
$\theta$ is the body pitch, $l$ is the variable leg length, and $\phi_{R,L}$ are the
right/left wheel rotations.

Three Pfaffian (non-holonomic) rolling constraints $A(\bq)\bdq = 0$ enforce pure rolling
without slipping:
\begin{align}
    -\dot{x}\sin\psi + \dot{y}\cos\psi &= 0 \quad \text{(no lateral slip)} \\
    (\dot{x}\cos\psi + \dot{y}\sin\psi) + \tfrac{d}{2}\dot\psi - r\dot\phi_R &= 0 \\
    (\dot{x}\cos\psi + \dot{y}\sin\psi) - \tfrac{d}{2}\dot\psi - r\dot\phi_L &= 0
\end{align}
These reduce the system to $m = 4$ independent pseudo-velocities
$\bm{u} = [\dot\phi_R, \dot\phi_L, \dot\theta, \dot{l}]^T$, related to $\bdq$ via the
null-space basis $S(\bq) \in \Real^{7\times 4}$ such that $\bdq = S(\bq)\bm{u}$:
\begin{equation}
    S(\bq) = \begin{bmatrix}
        \tfrac{r}{2}\cos\psi & \tfrac{r}{2}\cos\psi & 0 & 0 \\[2pt]
        \tfrac{r}{2}\sin\psi & \tfrac{r}{2}\sin\psi & 0 & 0 \\[2pt]
        \tfrac{r}{d} & -\tfrac{r}{d} & 0 & 0 \\[2pt]
        0 & 0 & 1 & 0 \\
        0 & 0 & 0 & 1 \\
        1 & 0 & 0 & 0 \\
        0 & 1 & 0 & 0
    \end{bmatrix}
    \label{eq:Smatrix}
\end{equation}
The forward velocity and yaw rate are kinematically determined:
$v_f = \tfrac{r}{2}(\dot\phi_R + \dot\phi_L)$,\quad
$\dot\psi = \tfrac{r}{d}(\dot\phi_R - \dot\phi_L)$.

Applying the Lagrange--d'Alembert projection $S^T(\bM\dot{\bm{u}} + \bm{n}) = S^T\bm{F}$
yields the \emph{reduced} equations:
\begin{equation}
    \bM_b(\bq)\dot{\bm{u}} + \bm{n}_b(\bq,\bm{u}) = \bB_b\btau
\end{equation}
with the input mapping matrix:
\begin{equation}
    \bB_b = \begin{bmatrix}
        1 & 0 & 0 \\
        0 & 1 & 0 \\
        -1 & -1 & 0 \\
        0 & 0 & 1
    \end{bmatrix}
    \label{eq:Bb}
\end{equation}
where $\btau = [\tau_R,\,\tau_L,\,F_l]^T$.  The critical observation is that the pitch row
of $\bB_b$ is $[-1,\,-1,\,0]$: the pitch $\theta$ is controlled \emph{only} by the
opposite reaction of the sum of wheel torques $-(\tau_R + \tau_L)$.  This confirms the
single degree of underactuation.

\subsection{Actuation and Underactuation}

Three independent actuators are available: wheel torque $\tau_w$ (or equivalently $\tau_R$,
$\tau_L$), knee torque $\tau_1$, and hip torque $\tau_2$.  The torso angle $q_3$ is
\textbf{unactuated} and must be controlled indirectly via centroidal coupling.

The ankle angle satisfies the algebraic constraint $q_0 = q_1 - q_2 + q_3$ and is not an
independent state.


% ============================================================
\section{Dynamic Model}
\label{sec:model}

\subsection{Lagrangian Derivation}

The absolute barycentre positions (wheel--ground contact as origin) are:
\begin{align}
    x_1 &= s + \tfrac{l_1}{2}\sin\phi_s, \quad
    y_1 = R_w + \tfrac{l_1}{2}\cos\phi_s \\[3pt]
    x_2 &= s + l_1\sin\phi_s + \tfrac{l_2}{2}\sin\phi_t, \quad
    y_2 = R_w + l_1\cos\phi_s + \tfrac{l_2}{2}\cos\phi_t \\[3pt]
    x_3 &= s + l_1\sin\phi_s + l_2\sin\phi_t + \tfrac{l_3}{2}\sin q_3, \quad
    y_3 = R_w + l_1\cos\phi_s + l_2\cos\phi_t + \tfrac{l_3}{2}\cos q_3
\end{align}
where $\phi_s = q_1 - q_2 + q_3$ (absolute shank angle) and $\phi_t = q_3 - q_2$ (absolute
thigh angle).

The kinetic energy is:
\begin{equation}
    T = \tfrac{1}{2}(m_w + I_w/R_w^2)\dot{s}^2 + \sum_{k=1}^{3}\!\left[\tfrac{1}{2}m_k
    \|\dot{\bm{r}}_k\|^2 + \tfrac{1}{2}I_k\dot{\phi}_k^2\right]
\end{equation}
and the potential energy: $V = \sum_k m_k g y_k$.

The Euler--Lagrange equations yield the standard structured form:
\begin{equation}
    \bM(\bq)\bddq + \bC(\bq,\bdq)\bdq + \bG(\bq) = \bB\btau
    \label{eq:eom}
\end{equation}
with the actuation matrix:
\begin{equation}
    \bB = \begin{bmatrix}
        -1 & 1 & 0 \\  1 & 0 & 0 \\  1 & 0 & 1 \\  -1 & 0 & 0
    \end{bmatrix} \in \Real^{4\times3}
\end{equation}

\subsection{Properties of the Dynamic Matrices}

The mass matrix $\bM(\bq)$ is symmetric and positive definite for all non-singular
configurations.  An important structural property is the \emph{skew-symmetry} of
$\dot{\bM} - 2\bC$: for any vector $\bm{v}$,
$\bm{v}^T(\dot{\bM} - 2\bC)\bm{v} = 0$.  This property is exploited by the SMC to
guarantee that the time derivative of the Lyapunov function along the system trajectories
contains only the sliding-surface error term, without parasitic energy-like cross-terms.

\subsection{Equilibrium and Static Verification}

At the upright equilibrium ($\bq = \bm{0}$, $\bdq = \bm{0}$), the gravitational vector
$\bG|_{\bq=0} = \bm{0}$, verified numerically.  For the PID controller, the equilibrium
posture at a desired CoM height $h_{des} = 0.75$~m is found by solving three simultaneous
equations with a numerical root-finder:
\begin{enumerate}
    \item \emph{Balance}: the mass-weighted horizontal CoM coordinate equals zero.
    \item \emph{Posture}: $q_1 + q_2 + q_3 = \pi/2$ (torso vertical).
    \item \emph{Height}: $l_1\sin q_1 + l_2\sin(q_1+q_2) + l_3\sin(q_1+q_2+q_3) = h_{des}$.
\end{enumerate}

\subsection{Decoupled VL-WIP Model}
\label{sec:vlwip}

For the LQR and MPC design, the multi-link leg is replaced by an \emph{equivalent centroid}:
a point mass $m_b$ at distance $l$ from the wheel axle, tilted at angle $\theta$ from the
vertical.  The equivalent centroid parameters $(l,\theta,I_y,I_z)$ are computed from the
full kinematic chain at each leg configuration.

The linearised VL-WIP state-space model has state
$\bm{x} = [s,\theta,\phi_w,\dot{s},\dot{\theta},\dot{\phi}_w]^T$ and input
$\bm{u} = [\tau_L,\tau_R]$.  The non-zero coefficients of the $(A,B)$ matrices
(eqs.~14 of~\cite{cui2022}) are:
\begin{align}
    a_1 &= \frac{-g l^2 m_b^2 r^2}{\Delta}, \qquad
    a_2 = \frac{g l m_b (2I_w + (m_b + 2m_w)r^2)}{\Delta} \\[4pt]
    b_1 &= \frac{r(I_y + l m_b(l+r))}{\Delta}, \qquad
    b_2 = \frac{-(2I_w + r(l m_b + (m_b + 2m_w)r))}{\Delta} \\[4pt]
    b_3 &= \frac{d\,r}{2 I_z r^2 + d^2(I_w + m_w r^2)}
\end{align}
where $\Delta = 2 I_w(I_y + m_b l^2) + (2 l^2 m_b m_w + I_y(m_b + 2m_w))r^2$.

The coefficient $b_3$ governs the yaw dynamics and is independent of pitch, while
$a_1,a_2,b_1,b_2$ vary with leg length $l$ and inertia $I_y$, motivating the
time-varying gain schedule.

\subsection{Upper-Body MPC Model}

For the MPC, the lumped-mass upper body is further decoupled into two independent
subsystems (eq.~9 of~\cite{cui2022}):

\textbf{Horizontal:} state $[s,\dot{s}]$, input $\Delta s$ (CoM offset).  The coupling is
$\ddot{s} = (g + \ddot{z}_{ref})/h_0 \cdot \Delta s$, essentially a virtual LIPM
(Linear Inverted Pendulum Model).

\textbf{Vertical:} state $[z,\dot{z}]$, input $F_z$ (vertical force).  The dynamics are
$\ddot{z} = F_z/m_b - g$.

Combined as a 5-state system $[s,\dot{s},z,\dot{z},\Delta s]$ with inputs
$[\dot{\Delta s}, F_z]$:
\begin{equation}
    A_{MPC} = \begin{bmatrix}
        0 & 1 & 0 & 0 & 0 \\
        0 & 0 & 0 & 0 & g/h_0 \\
        0 & 0 & 0 & 1 & 0 \\
        0 & 0 & 0 & 0 & 0 \\
        0 & 0 & 0 & 0 & 0
    \end{bmatrix},\quad
    B_{MPC} = \begin{bmatrix}
        0 & 0 \\ 0 & 0 \\ 0 & 0 \\ 0 & 1/m_b \\ 1 & 0
    \end{bmatrix}
\end{equation}
with nominal height $h_0 = 0.6$~m and upper-body mass $m_b = 66.5$~kg.

\subsection{CoM Kinematics and Pitch Jacobian}

The horizontal and vertical CoM offsets relative to the wheel axle are:
\begin{align}
    S_C &= C_a\sin\phi_s + C_b\sin\phi_t + C_c\sin q_3 \\
    Z_C &= C_a\cos\phi_s + C_b\cos\phi_t + C_c\cos q_3
\end{align}
with $C_a = (m_1 l_1/2 + (m_2+m_3)l_1)/m_b$, $C_b = (m_2 l_2/2 + m_3 l_2)/m_b$,
$C_c = m_3 l_3/(2m_b)$.  The upper-body pitch angle and its analytic Jacobian are:
\begin{equation}
    \theta_b = \arctan2(S_C,Z_C), \qquad
    J_\theta = \frac{Z_C J_S - S_C J_Z}{S_C^2 + Z_C^2}
\end{equation}
The forward CoM velocity is $v_{CoM} = R_w\dot{q}_w + \kappa\,\dot{S}_C$, where
$\kappa = m_b/m_t$ is the body-to-total mass ratio.  The analytic Jacobians $J_S$, $J_Z$
and the bias terms $h_\theta$, $h_v$ (Lie-derivative corrections accounting for
Coriolis/centripetal accelerations in the output space) are computed at every timestep,
eliminating the need for numerical differentiation.


% ============================================================
\section{Cascade PID Controller}
\label{sec:pid}

\subsection{Architecture}

The simplest controller is a two-level cascade PID, inspired by Fig.~6 of~\cite{cui2022}.

\textbf{Outer loop (velocity PI).}  Maps wheel velocity error
$e_v = \dot{q}_w^* - \dot{q}_w$ to a shank-angle offset:
\begin{equation}
    \delta q_1 = -(K_{P,v}\,e_v + K_{I,v}\!\int e_v\,dt),
    \quad K_{P,v} = 0.08,\ K_{I,v} = 0.06
\end{equation}
with saturation $|\delta q_1| \leq 0.2$~rad to prevent extreme lean demands.  The negative
sign reflects the non-minimum phase nature of the plant.

\textbf{Inner loop (balance PD).}  Stabilises the shank angle to $q_1^* + \delta q_1$:
\begin{equation}
    \tau_w = K_{P,w}(q_1^* + \delta q_1 - q_1) + K_{D,w}(0 - \dot{q}_1),
    \quad K_{P,w} = 400,\ K_{D,w} = 80\ \text{[Nm/rad]}
\end{equation}

\textbf{Posture PD.}  Two independent loops maintain knee and hip at the equilibrium
targets (computed as described in Section~\ref{sec:model}):
\begin{equation}
    \tau_k = K_{P,j}(q_k^* - q_k) + K_{D,j}(0 - \dot{q}_k),
    \quad K_{P,j} = 800,\ K_{D,j} = 50\ \text{[Nm/rad]}
\end{equation}

The simulation starts with an initial perturbation of $-0.1$~rad on $q_1$ from equilibrium
and is integrated over 20~s using ODE45 with the full 9-state dynamics (4~positions + 4~velocities + 1~integral state).

\subsection{Design Choices and Failure Analysis}

The gains are tuned empirically without model knowledge.  The cascade structure assumes that
the inner loop (balance) is significantly faster than the outer loop (velocity), but the
actual plant dynamics violate this assumption: the coupling between pitch and posture means
that the posture PD generates torques that perturb the pitch faster than the balance PD can
compensate.

The progressive failure mechanism is as follows.  A small initial lean causes the velocity
PI to demand a compensating shank offset; however, the high posture gains
($K_{P,j} = 800$) produce a counter-torque that pushes the pitch in the opposite direction.
The resulting oscillation grows because the wheel torque saturates at $\pm 6$~Nm, leaving
the balance loop open for the portion of the cycle where $|\tau_w| = 6$~Nm.  This
saturation-induced instability is a well-known limitation of linear controllers applied to
underactuated systems with bounded inputs.


% ============================================================
\section{LQR + MPC Controller}
\label{sec:lqr_mpc}

\subsection{TV-LQR Balance Loop}

The linearised VL-WIP model (Section~\ref{sec:vlwip}) is scheduled as a function of leg
length $l$ via a Time-Varying LQR (TV-LQR).  At construction time, 25~hip-angle samples
are taken from $[-0.45, +0.45]$~rad (with knee $= -2\cdot\text{hip}$ to keep the axle
under the hip).  For each sample, the equivalent centroid provides $(l, I_y)$, from which
the VL-WIP $(A,B)$ matrices are built.

The full-state gain $K(l)$ is obtained by solving the \emph{Continuous Algebraic Riccati
Equation} (CARE) at each grid point.  At runtime, $K$ is linearly interpolated on the
$l$-grid.  The CARE is solved via the Hamiltonian eigendecomposition:
\begin{equation}
    H = \begin{bmatrix} A & -B R^{-1} B^T \\ -Q & -A^T \end{bmatrix}
\end{equation}
The $n$ eigenvectors of $H$ with negative real-part eigenvalues form the stable subspace
$[V_{11};V_{21}]$; then $P = \text{Re}(V_{21}V_{11}^{-1})$ and $K = R^{-1}B^T P$.

\subsubsection{State and Input Cost Matrices}

The state cost matrix:
\begin{equation}
    Q = \operatorname{diag}(30,\ 400,\ 80,\ 15,\ 6,\ 2)
\end{equation}
penalises pitch most heavily ($Q_{22} = 400$) for rapid tilt recovery, while position
($Q_{11} = 30$) maintains station-keeping without fighting the balance loop.  The yaw
penalty ($Q_{33} = 80$) provides heading hold.

The input cost is constructed in \emph{common-mode / differential-mode} coordinates:
\begin{equation}
    R = T_{cd}^T \operatorname{diag}(R_{cm}, R_{diff}) T_{cd}, \quad
    T_{cd} = \begin{bmatrix} 0.5 & 0.5 \\ -0.5 & 0.5 \end{bmatrix}
\end{equation}
with $R_{cm} = 2.0$ and $R_{diff} = 100.0$.  The 50$\times$ higher differential penalty
keeps the yaw loop crossover ($\sim$16~rad/s) well below the torsional resonance of the
stance legs ($\sim$76~rad/s $\approx 12$~Hz: torso yaw inertia on the VMC springs), which
a high yaw gain would excite into a permanent limit cycle.

\subsection{Upper-Body MPC}

A linear MPC plans the CoM offset $\Delta s$ and vertical force $F_z$ over a finite
horizon.

\subsubsection{QP Formulation}

The continuous-time model is discretised at $\Delta t_{MPC} = 0.02$~s.  The condensed QP
eliminates the state trajectory $x_i = \Phi_i x_0 + \Gamma_i U + C_i$:
\begin{equation}
    \min_U \; \tfrac{1}{2}U^T H\,U + g^T U \quad \text{s.t.} \quad
    \bm{l} \leq U \leq \bm{h}
\end{equation}
where $H = \Gamma^T \bar{S}\,\Gamma + W\cdot I_N$ and
$g = \Gamma^T \bar{S}(\Phi x_0 + C - x_{ref})$, with $\bar{S}$ the block-diagonal stage
cost.

\textbf{Parameters:}
\begin{itemize}[leftmargin=*]
    \item Horizon: $N = 25$ steps $\times$ 0.02~s $= 0.5$~s preview.
    \item Execution rate: 500~Hz / 5 $= 100$~Hz.
    \item Horizontal stage cost: $S_h = \operatorname{diag}(50, 20)$ on $(s, \dot{s})$;
          input weight $W_h = 200$ on $\Delta s$.
    \item Vertical stage cost: $S_v = \operatorname{diag}(5000, 150)$ on $(z, \dot{z})$;
          input weight $W_v = 0.002$ on $F_z$.
    \item $\Delta s$ bounds: $|\Delta s| \leq \min(\mu z,
          \sqrt{L_{max}^2 - z_b^2}, 0.03)$~m (friction cone, geometric reach, hard cap).
    \item $F_z$ bounds: $[0.3,\ 3.0] \times m_b g$ (normal stance).
\end{itemize}

\subsubsection{FISTA Solver}

The box-constrained QP is solved with the FISTA (Fast Iterative Shrinkage-Thresholding
Algorithm) --- an accelerated projected gradient method:
\begin{align}
    x_{k+1} &= \Pi_{[\bm{l},\bm{h}]}(y_k - \tfrac{1}{L}(Hy_k + g)) \\
    t_{k+1} &= \tfrac{1}{2}(1 + \sqrt{1 + 4t_k^2}) \\
    y_{k+1} &= x_{k+1} + \tfrac{t_k - 1}{t_{k+1}}(x_{k+1} - x_k)
\end{align}
with step size $1/L$ where $L$ is the largest eigenvalue of $H$, and 60~iterations per
solve.  The solver is warm-started with the shifted previous solution (drop first element,
duplicate last) for temporal coherence.

\subsubsection{MPC--LQR Integration}

Between MPC solves (at 500~Hz), the LQR reference $s_{plan}$ follows the MPC-predicted CoM
trajectory via linear interpolation over one MPC step.  The LQR state reference becomes:
\begin{equation}
    X_{ref} = [s_{plan} - \Delta s,\;\theta_{ref},\;\psi_{ref},\;
    \dot{s}_{plan},\;0,\;\dot\psi_{ref}]^T
\end{equation}
where $\theta_{ref} = \arctan(\Delta s / z_{ref})$ maps the MPC's horizontal offset into a
pitch-angle setpoint.

\subsection{VMC Leg Controller}

The MPC outputs $(\Delta s, z_{ref}, F_z)$ are converted to hip/knee torques via a
task-space impedance controller (Virtual Model Controller).  Each leg is an independent
planar 2-link chain controlled in Cartesian $(x,z)$:
\begin{equation}
    \bm{F} = K_p(\bm{p}_d - \bm{p}_f) + K_d(\bm{v}_d - \bm{v}_f) + \bm{F}_{ff},
    \qquad \btau_{leg} = J^T\bm{F}
\end{equation}

The gains are phase-dependent (Table~\ref{tab:vmc_gains}).

\begin{table}[H]
    \centering
    \caption{VMC gains per phase (per leg, base-frame $x$-$z$).}
    \label{tab:vmc_gains}
    \begin{tabular}{lcc}
        \toprule
        Phase & $K_p$ [N/m] & $K_d$ [Ns/m] \\
        \midrule
        Hold (pre-release) & diag(1500, 2000) & diag(40, 50) \\
        Stance (MPC active) & diag(1500, 0) & diag(40, 15) \\
        Flight & diag(800, 800) & diag(20, 20) \\
        \bottomrule
    \end{tabular}
\end{table}

A key design choice is that in stance with MPC, the \emph{vertical stiffness is zero}
($K_{p,zz} = 0$): the vertical support force comes entirely from the MPC feedforward
$F_z$, while only damping (15~Ns/m) acts vertically.  This ensures the MPC retains full
authority over the CoM height.  The feedforward force per leg is
$\bm{F}_{ff} = -\tfrac{1}{2}F_z\,\hat{\bm{g}}_b$, where $\hat{\bm{g}}_b$ is the gravity
unit vector rotated into the body frame.

\subsection{Yaw Controller}

The yaw is controlled by a PD law with Coulomb/viscous friction feedforward and a slow
integral:
\begin{equation}
    \tau_{diff} = K_{P,\psi} e_\psi + K_{D,\psi}\dot{e}_\psi +
    K_{I,\psi}\!\int_{\!|e_\psi|>\delta_{db}} e_\psi\,dt +
    \mu_C\tanh(\dot\psi_{ref}/0.01) + \mu_v\dot\psi_{ref}
\end{equation}
with $K_{P,\psi} = 0.89$, $K_{D,\psi} = 0.21$, $K_{I,\psi} = 0.6$~Nm/wheel,
$\mu_C = 0.25$~Nm/wheel (Coulomb, identified in Gazebo), $\mu_v = 0.10$~Nm$\cdot$s/wheel
(viscous), and deadband $\delta_{db} = 0.035$~rad to avoid stick-slip hunting.

This controller is \emph{deliberately not} sliding mode: a discontinuous term here would
excite the 12.5~Hz torsional resonance of the stance legs into a permanent limit cycle.
The PD poles reproduce the closed-loop poles of the TV-LQR ($-7.45 \pm 2.62j$).

\subsection{Torque Priority and Traction Limits}

The wheel torques are decomposed into common mode (balance) and differential mode (yaw).
The differential torque is clamped first: $|\tau_{diff}| \leq 2.5$~Nm/wheel.  The
remaining traction budget is assigned to balance:
\begin{equation}
    |\tau_{cm}| \leq \max\!\left(\min(10, \mu_{safe}\cdot\mu\cdot N_{wheel}\cdot r) -
    |\tau_{diff}|,\;0\right)
\end{equation}
where $\mu_{safe} = 0.8$ is a friction safety factor and $N_{wheel}$ is the estimated
per-wheel normal force (floored at 35\% of body weight to prevent traction budget
collapse during transients).

\subsection{Kalman State Estimator}

Three independent scalar Kalman filters (one per world axis $x,y,z$) run at 500~Hz.  Each
has state $[p,v]^T$ with a constant-acceleration process model:
\begin{equation}
    F = \begin{bmatrix} 1 & dt \\ 0 & 1 \end{bmatrix},\quad
    B_u = \begin{bmatrix} \tfrac{1}{2}dt^2 \\ dt \end{bmatrix},\quad
    Q_k = q_{acc}\,B_u B_u^T,\quad
    R_k = \operatorname{diag}(10^{-4}, 10^{-3})
\end{equation}
with $q_{acc} = 0.5$.  The input is the gravity-compensated, world-frame IMU acceleration.
\textbf{Prediction} runs every tick; \textbf{correction} (stance only) uses wheel odometry
for position/velocity, offset by leg forward kinematics (axle-to-body transform).  In
flight, the filter runs in pure prediction mode.


% ============================================================
\section{Sliding Mode Controller}
\label{sec:smc}

\subsection{Architecture Overview}

The SMC is a nonlinear cascade with two time-scales:
\begin{enumerate}
    \item \textbf{Outer loop (slow, $\sim$50~Hz bandwidth):} PI on CoM velocity error
          $\to$ pitch reference $\theta^*$.
    \item \textbf{Inner loop (fast, 500~Hz):} three SMC channels on the output vector
          $\bm{y} = [\theta_b,\,q_1,\,q_2]^T$.
\end{enumerate}

\subsection{Non-Minimum Phase Outer Loop}

The outer loop maps velocity error to pitch reference:
\begin{equation}
    \theta^* = \operatorname{sat}\!\left(K_P e_v + K_I\!\int\!e_v\,dt,\;
    \pm\theta_{max}\right)
\end{equation}
with $K_P = -0.30$ (ROS~2) or $-0.22$ (MATLAB), $K_I = -0.10$ (resp.\ $-0.06$),
$\theta_{max} = 0.21$~rad ($\approx 12\dgr$).

The gains are \textbf{negative} because the system is non-minimum phase: to accelerate
forward, the robot must first lean forward (moving the axle the ``wrong'' way).  The
position error is folded into the velocity reference as
$v_{ref} = v_d + K_{pos}(s_{ref} - s)$ with $K_{pos} = 1.2$~s$^{-1}$ and
$|v_{ref}| \leq 2.0$~m/s.

\textbf{Anti-windup.}  The integrator $I_v$ is updated only when the raw pitch reference
has not hit the saturation limit: $\dot{I}_v = e_v$ if $|\theta_{raw}^* - \theta^*| <
\epsilon$, else $\dot{I}_v = 0$.  Both integrators ($I_v$ and the twisting integrator $W$)
are zeroed at every phase transition to avoid torque discontinuities.

\subsection{Super-Twisting Pitch Channel}

Channel~1 uses the super-twisting second-order SMC~\cite{levant1993, moreno2012}, which
provides a \emph{continuous} control signal:
\begin{align}
    s_1 &= \dot{\theta}_b + c_1(\theta_b - \theta^*) \\[4pt]
    \nu_1 &= -k_a\sqrt{|s_1|}\,\frac{s_1}{|s_1|+\varepsilon} + W \\[4pt]
    \dot{W} &= -k_b\,\frac{s_1}{|s_1|+\varepsilon}
\end{align}

The regularised sign function $s/(|s|+\varepsilon)$ replaces the ideal
$\operatorname{sgn}(s)$, eliminating the discontinuity at $s=0$ and the associated
chattering.  The parameters differ between the MATLAB prototype and the ROS~2
implementation (Table~\ref{tab:smc_params}), reflecting the different model fidelities.

\begin{table}[H]
    \centering
    \caption{Super-twisting parameters: MATLAB (planar, no noise) vs.\ ROS~2
    (3D, sensor noise).}
    \label{tab:smc_params}
    \begin{tabular}{lcc}
        \toprule
        Parameter & MATLAB & ROS~2 \\
        \midrule
        $c_1$ [s$^{-1}$] & 6.0 & 10.0 \\
        $k_a$ & 3.5 & 25.0 \\
        $k_b$ & 12.0 & 250.0 \\
        $\varepsilon$ & 0.01 & 0.02 \\
        $|W|_{max}$ [rad/s$^2$] & --- & 60.0 \\
        $\tau_{wheel,max}$ [Nm] & 25.0 & 10.0 \\
        \bottomrule
    \end{tabular}
\end{table}

The ROS~2 gains are significantly more aggressive ($k_a = 25$, $k_b = 250$) because the
higher-fidelity simulation includes sensor noise, actuator lag, and contact dynamics that
require faster convergence.  The design comment notes: ``$c_1$ multiplies $\dot\theta$, a
filtered numerical derivative ($\tau \approx 9$~ms).  On the linearised plant the loop
survives $c_1 = 40$; what grows with $c_1$ is torque activity, not instability.  The real
limit is the 12.5~Hz torsional resonance of the stance legs.''

\textbf{Anti-windup on the twisting integrator.}  In MATLAB, $\dot{W} = 0$ whenever
\emph{any} of the three torques saturates.  In ROS~2, $\dot{W} = 0$ only when the common
wheel torque saturates (more permissive, since the leg torques saturating does not affect
the pitch channel).  Additionally, the Euler integration step is capped:
$dt_{eff} = \min(dt, 4\text{~ms})$ to prevent a late control tick from integrating $k_b$
over an unreasonably long interval.

\subsection{Posture Channels (Knee and Hip)}

Channels~2 and~3 use classical first-order SMC with boundary layer:
\begin{equation}
    \nu_k = -k_k\operatorname{sat}\!\left(\frac{s_k}{\phi_k}\right) - \eta_k s_k,
    \quad k = 2,3
\end{equation}
with $c_{2,3} = 12$~s$^{-1}$, $k_{2,3} = 60$, $\eta_{2,3} = 20$,
$\phi_{2,3} = 0.05$~rad.  The boundary layer of width $\phi_k$ converts the discontinuous
sign function into a continuous saturation, while $\eta_k s_k$ provides exponential
attraction within the boundary layer.

\subsection{Torque Inversion via Partial Feedback Linearisation}

\subsubsection{MATLAB Implementation}

The three virtual inputs $\bm{\nu}$ are mapped to actual torques by inverting the
output dynamics:
\begin{equation}
    A_d = J_y\,\bM^{-1}\bB \in \Real^{3\times3}, \qquad
    b_d = -J_y\,\bM^{-1}(\bC\bdq + \bG) + h_y
\end{equation}
\begin{equation}
    \btau = A_d^{-1}(-b_d - c \odot \dot{\bm{e}} + \bm{\nu})
\end{equation}
where $J_y = [J_{\theta_b};\,\bm{e}_1^T;\,\bm{e}_3^T]$ is the 3$\times$4 output Jacobian
and $h_y$ collects the Lie-derivative correction terms.

\subsubsection{ROS~2 Implementation}

In the ROS~2 node, the pitch channel uses the \emph{VL-WIP linearised inversion} instead
of the full model inversion:
\begin{equation}
    \tau_{cm} = \frac{\nu_1 - a_2\theta}{2b_2}
\end{equation}
where $a_2, b_2$ are the VL-WIP coefficients (Section~\ref{sec:vlwip}) interpolated at
the current pendulum length.  The leg channels use per-leg task-space SMC (see below).

\subsection{Per-Leg Task-Space SMC (ROS~2)}

In the ROS~2 implementation, the posture is controlled by a per-leg task-space sliding mode
controller operating in Cartesian $(x,z)$:
\begin{equation}
    \bm{s} = (\bm{v}_f - \bm{v}_d) + c \odot (\bm{p}_f - \bm{p}_d), \qquad
    \bm{F} = -K_{lin}\,\bm{s} - K_{sw}\operatorname{clip}(\bm{s}/\phi,\,-1,\,1) +
    \bm{F}_{ff}
\end{equation}
with phase-dependent gains $c = [37.5,\,48.0]$~s$^{-1}$, $K_{lin} = \operatorname{diag}(40,\,25)$~N/(m/s),
$K_{sw} = 8.0$~N, $\phi = 0.03$~m/s.  The rationale is that the old VMC was already a
sliding-mode law without the switching term:
$K_p\bm{e} + K_d\dot{\bm{e}} \equiv -K_d[\dot{\bm{e}} + (K_p/K_d)\bm{e}]$, a surface
with $c = K_p/K_d$.  The parameters $c$ and $K_{lin}$ reproduce the validated VMC gains
exactly; only the robust switching term $K_{sw}$ is new.

\subsection{Singularity Monitoring}

A runtime diagnostic $\sigma = M_{22} + M_{42}$ monitors the effective inertia coupling
between the wheel torque and the pitch dynamics.  As $\sigma \to 0$, the wheel loses
authority over pitch (fully extended leg singularity).  The comment notes:
``$\sigma > 0$ always in operation; $\sigma \to 0$ means loss of wheel authority.''


% ============================================================
\section{Results and Discussion}
\label{sec:results}

All controllers are benchmarked in two scenarios: (i)~\textbf{free balance} (released from
rest, $v^* = 0$, 26~s simulation) and (ii)~\textbf{push recovery} (impulsive horizontal
force at $t = 4$~s).

\subsection{Free Balance Benchmark}

Figure~\ref{fig:dashboard} shows the overall dashboard.

\begin{figure}[H]
    \centering
    \includegraphics[width=\linewidth]{figures/dashboard_overview.png}
    \caption{Benchmark overview: pitch angle, wheel torque command, linear velocity, and
    contact normal force for PID (red), MPC/LQR (blue), and SMC (green).}
    \label{fig:dashboard}
\end{figure}

The PID controller (red) fails to stabilise the robot.  The pitch grows into a sustained
oscillation of amplitude $\pm 15\dgr$--$18\dgr$, the wheel torque saturates at $\pm 6$~Nm,
and the contact force fluctuates between 0~N and 42~N, indicating intermittent wheel
lift-off.  The velocity drifts sinusoidally with peaks of $\pm 3$~m/s.  This confirms the
saturation-induced instability analysed in Section~\ref{sec:pid}: once the wheel torque
rails, the balance loop is effectively open, and the pitch diverges until gravity pulls it
back through the opposite rail.

In contrast, both MPC (blue) and SMC (green) stabilise the robot within $\approx 2$~s and
thereafter maintain pitch within $\pm 0.5\dgr$, with wheel torque below 0.15~Nm --- three
orders of magnitude smaller than the PID effort.

\subsubsection{Position and Velocity Tracking}

Figure~\ref{fig:posvel} confirms that under MPC and SMC the robot stays stationary
($|x| < 0.02$~m, $|\dot{x}| < 0.05$~m/s), while the PID produces sustained drift of
amplitude $\pm 0.35$~m.  The drift is sinusoidal and phase-locked with the pitch
oscillation, confirming the coupling between the uncontrolled pitch and the position.

\begin{figure}[H]
    \centering
    \includegraphics[width=\linewidth]{figures/position_and_speed.png}
    \caption{Horizontal position $x$ and velocity $\dot{x}$ for the three controllers.}
    \label{fig:posvel}
\end{figure}

\subsubsection{Pitch and Torque Comparison}

Figure~\ref{fig:pitchtorque} isolates the pitch and torque channels.  The MPC achieves a
smoother pitch recovery thanks to the 0.5~s preview horizon of the MPC, which anticipates
the CoM trajectory and pre-positions the legs; the SMC is faster in the reaching phase
(finite-time convergence guaranteed by the super-twisting algorithm) but shows a residual
chattering of $\approx \pm 0.1\dgr$ from the boundary-layer approximation and the 500~Hz
discretisation.

The MPC torque signal is notably smoother: the TV-LQR produces a continuous optimal
feedback, whereas the SMC torque exhibits small-amplitude high-frequency oscillations from
the regularised sign function.  In a physical implementation, the MPC's smooth torque
profile would be less likely to excite structural resonances in the drivetrain.

\begin{figure}[H]
    \centering
    \includegraphics[width=\linewidth]{figures/pitch_and_torque.png}
    \caption{Pitch angle and wheel torque command.  The MPC provides the smoothest torque
    profile; the SMC shows small chattering from the boundary layer.}
    \label{fig:pitchtorque}
\end{figure}

\subsubsection{Full State Tracking and Errors}

Figure~\ref{fig:tracking_balance} provides the full state-vs-reference comparison.  The MPC
maintains $z = 0.220$~m with sub-millimetre error because the vertical DOF is directly
driven by the MPC's $F_z$ feedforward with the VMC's vertical stiffness set to zero; the
SMC settles at $z \approx 0.222$~m due to a residual offset from the fixed posture
set-points and nonlinear kinematics; the PID oscillates by $\pm 30$~mm.

\begin{figure}[H]
    \centering
    \includegraphics[width=\linewidth]{figures/tracking_and_errors.png}
    \caption{State vs desired reference (left) and tracking errors (right): pitch,
    position, velocity, and CoM height.}
    \label{fig:tracking_balance}
\end{figure}

\subsubsection{CoM Height and Ground Contact Force}

Figure~\ref{fig:comheight} details the CoM height regulation and contact force.  The SMC
height ripple ($\approx 3$~mm) originates from chattering in the posture channels
(boundary-layer width $\phi = 0.05$~rad generates a residual oscillation that propagates
through the leg kinematics to the CoM height).  The PID causes repeated wheel lift-off:
the contact force drops to 0~N whenever the pitch exceeds $\sim 12\dgr$, confirming that
the robot is on the verge of toppling.

\begin{figure}[H]
    \centering
    \includegraphics[width=\linewidth]{figures/com_height_and_forces.png}
    \caption{CoM height $z$ (top) and ground contact normal force $F_n$ (bottom).  PID
    causes wheel lift-off; MPC and SMC maintain continuous contact.}
    \label{fig:comheight}
\end{figure}

\subsection{Push Recovery Benchmark}

The PID is excluded (already unstable).  Figure~\ref{fig:push_tracking} shows the state
trajectory after the push.

\begin{figure}[H]
    \centering
    \includegraphics[width=\linewidth]{figures/tracking_and_errors.png}
    \caption{State vs reference and tracking errors after an impulsive push at $t = 4$~s.
    Both MPC and SMC recover.}
    \label{fig:push_tracking}
\end{figure}

The MPC (blue) produces a compliant response: pitch peaks at $\approx 12\dgr$, recovery in
$\approx 3$~s via a smooth oscillatory transient shaped by the LQR cost trade-off between
fast convergence and low torque.  The SMC (green) reacts more aggressively: the
super-twisting drives $s_1$ back to the chattering band in $\approx 1.5$~s, exploiting the
full traction budget.  The faster recovery comes at the cost of higher peak torque and
more aggressive wheel acceleration.

\subsubsection{Phase Portrait and Disturbance Rejection}

Figure~\ref{fig:phase_push} shows the $(\theta,\dot\theta)$ phase portrait, which reveals
the fundamental difference between the two control strategies.

\begin{figure}[H]
    \centering
    \includegraphics[width=\linewidth]{figures/disturbance_and_phase_portrait.png}
    \caption{Disturbance rejection (left) and $(\theta,\dot\theta)$ phase portrait (right).
    Dashed purple: sliding manifold $s_1 = \dot\theta + c_1\theta = 0$.}
    \label{fig:phase_push}
\end{figure}

The MPC orbit follows a large, smooth arc that crosses the sliding manifold multiple times
--- its trajectory is dictated by the quadratic cost, not by the geometry of a switching
surface.  The SMC orbit, by contrast, immediately seeks the manifold
$s_1 = \dot\theta + c_1\theta = 0$ and then slides along it toward the origin, consistent
with second-order SMC theory.  The inward spiral confirms finite-time convergence of the
super-twisting algorithm.

\subsubsection{Super-Twisting Sliding Variable}

Figure~\ref{fig:smc_surface} provides direct evidence of the super-twisting algorithm's
behaviour.

\begin{figure}[H]
    \centering
    \includegraphics[width=\linewidth]{figures/smc_sliding_surface.png}
    \caption{Sliding variable $s_1(t)$ (left) and SMC state trajectory in
    $(\theta,\dot\theta)$ (right).}
    \label{fig:smc_surface}
\end{figure}

Before the push ($t < 4$~s), $|s_1| < 0.1$~rad/s (chattering band).  At impact, $s_1$
drops sharply to $-3.5$~rad/s.  The super-twisting algorithm then drives $s_1$ back to
the chattering band in $\approx 1.5$~s, with the characteristic $\sqrt{|s|}$ convergence
rate.  The right panel shows the state spiralling inward toward the sliding manifold.  The
integral state $W$ absorbs the persistent disturbance (gravity coupling), keeping $s_1$
centred near zero once the reaching phase is complete.

\subsection{Quantitative Comparison}

Table~\ref{tab:comparison} summarises the key performance metrics.

\begin{table}[H]
    \centering
    \caption{Performance summary --- free balance scenario (26~s simulation).}
    \label{tab:comparison}
    \begin{tabular}{lccc}
        \toprule
        Metric & PID & LQR+MPC & SMC \\
        \midrule
        Settling time (pitch $< 0.5\dgr$) & $\infty$ & $\approx 2.0$~s & $\approx 1.8$~s \\
        Steady-state pitch RMS & N/A & $< 0.2\dgr$ & $\approx 0.3\dgr$ \\
        Max position drift $|x|$ & $\pm$350~mm & $<$20~mm & $<$20~mm \\
        CoM height error RMS & $\sim$25~mm & $<$1~mm & $\approx$2~mm \\
        Wheel torque RMS (steady) & $\pm$6~Nm & $<$0.05~Nm & $<$0.15~Nm \\
        Chattering proxy $\overline{|\Delta\tau|}$ & N/A & $<$0.01~Nm & $\approx$0.08~Nm \\
        Contact loss events & Many & None & None \\
        Computational cost & Negligible & Medium (QP) & Low (analytic) \\
        Model knowledge & None & Full linear & Nonlinear \\
        Robustness guarantee & No & Quadratic & Finite-time \\
        \bottomrule
    \end{tabular}
\end{table}

\subsection{Discussion}

\paragraph{PID.}  The cascade PID is conceptually simple but lacks model knowledge.  The
wheel torque saturation triggers a limit cycle: when $|\tau_w| = \tau_{max}$, the balance
loop is open and the pitch diverges until gravity reverses the motion.  This is a
structural failure mode that no amount of gain tuning can fully eliminate for the given
actuator limits.  A gain schedule on leg length could delay but not prevent the instability.

\paragraph{LQR + MPC.}  The decoupled architecture separates concerns optimally: the
TV-LQR handles fast balance with guaranteed quadratic stability (via the Riccati solution),
the MPC previews the CoM trajectory over 0.5~s and pre-positions the legs.  The
common/differential decomposition of the wheel torque elegantly separates balance from yaw
without cross-coupling.  The smooth torque profile is the main practical advantage over SMC
for hardware deployment.  The limitation is the linear model assumption: for perturbations
beyond $\sim 15\dgr$ pitch, the VL-WIP linearisation degrades.

\paragraph{SMC.}  The super-twisting pitch channel provides provably finite-time
convergence to the sliding manifold regardless of bounded disturbances and model
uncertainties~\cite{moreno2012}.  The partial feedback linearisation ensures exact
decoupling at the operating point.  The residual chattering ($\approx 0.08$~Nm mean
torque variation per sample) is a consequence of: (i)~the 500~Hz discretisation of the
continuous-time super-twisting law, (ii)~the regularisation parameter $\varepsilon = 0.02$,
and (iii)~unmodelled dynamics (torsional resonance of the legs).  In hardware, joint
compliance and motor inductance would filter this residual.

The key design contribution is the \textbf{super-twisting} algorithm for the pitch channel:
it provides a continuous control signal while preserving sliding-mode robustness.  The
integrator freeze at saturation prevents the windup that would otherwise accumulate after a
large push --- a feature absent in classical first-order SMC designs.


% ============================================================
\section{Contributions with Respect to Class Topics}
\label{sec:contributions}

\begin{itemize}[leftmargin=*]
    \item \textbf{Wheeled robot kinematics and dynamics:} the rolling constraint is derived
          from first principles; the 7-DOF model is reduced to 4~DOF via null-space
          projection (Lagrange--d'Alembert); the VL-WIP equivalent centroid provides a
          configuration-dependent linearisation for the LQR.
    \item \textbf{Balance and underactuated control:} three controllers span the spectrum
          from purely reactive (PID) to model-based predictive (MPC) to provably robust
          nonlinear (SMC), with a principled quantitative comparison.
    \item \textbf{Sliding mode control:} the super-twisting algorithm, boundary-layer
          technique, and partial feedback linearisation are applied to a multi-DOF plant,
          with per-leg task-space SMC extending the approach to Cartesian space.
    \item \textbf{Optimal control (LQR/MPC):} TV-LQR scheduled on leg length via CARE;
          condensed QP with FISTA solver at 100~Hz; common/differential torque decomposition
          to separate balance from yaw.
    \item \textbf{State estimation:} three-axis linear Kalman filter fusing IMU, wheel
          odometry, and leg forward kinematics, with stance/flight mode switching.
    \item \textbf{Software-in-the-loop:} full pipeline from symbolic Lagrangian (MATLAB) to
          closed-loop ROS~2/Gazebo simulation with automated benchmarking, CSV logging, and
          publication-quality plot generation.
\end{itemize}


% ============================================================
\section{Conclusions}
\label{sec:conclusions}

This report presented the modelling and control of RoTino, a simulated VL-WIP wheeled biped
robot, covering the full pipeline from topology/DOF analysis and Lagrangian derivation to
multi-controller benchmarking in both MATLAB and ROS~2/Gazebo.

Key conclusions:
\begin{enumerate}
    \item The Lagrangian model --- derived symbolically and validated at equilibrium ---
          captures the essential nonlinear dynamics and runs in real time at 500~Hz in both
          MATLAB and Python.
    \item The cascade PID fails due to saturation-induced instability: the limited wheel
          torque opens the balance loop during large pitch excursions, creating a limit cycle
          that no fixed-gain linear controller can eliminate.
    \item Both LQR+MPC and SMC achieve stable balance and push recovery.  The MPC provides
          the smoothest torque profile and the most accurate CoM height tracking
          ($< 1$~mm error) thanks to the preview horizon and the zero-vertical-stiffness VMC
          design.  The SMC achieves faster reaching ($\sim 1.5$~s vs.\ 3~s for push
          recovery) and stronger robustness guarantees (finite-time convergence) at lower
          computational cost (analytic inversion vs.\ QP solve).
    \item The super-twisting second-order SMC is identified as the best practical choice for
          embedded real-time implementation: it is continuous, computationally inexpensive,
          and provides finite-time convergence guarantees independent of model accuracy.
    \item The two-stage development pipeline (MATLAB prototyping $\to$ ROS~2/Gazebo
          validation) proved effective: the MATLAB gains served as initialisation for the
          ROS~2 controllers, which were then re-tuned for the higher-fidelity simulation.
\end{enumerate}

Future work could extend the model to full 3D (coronal-plane dynamics, inertial coupling),
add joint compliance modelling to predict chattering attenuation, investigate
reinforcement-learning-based gain scheduling, and validate the controllers on physical
hardware.


% ============================================================
\begin{thebibliography}{9}

\bibitem{cui2022}
Z.\ Cui, Z.\ Zhao, H.\ Zhu, J.\ Bao, and Z.\ Zhang,
``Modeling and Control of a Wheeled Biped Robot,''
\emph{Micromachines}, vol.\ 13, no.\ 5, p.\ 747, 2022.

\bibitem{levant1993}
A.\ Levant,
``Sliding order and sliding accuracy in sliding mode control,''
\emph{Int.\ J.\ Control}, vol.\ 58, no.\ 6, pp.\ 1247--1263, 1993.

\bibitem{moreno2012}
J.\ A.\ Moreno and M.\ Osorio,
``Strict Lyapunov functions for the super-twisting algorithm,''
\emph{IEEE Trans.\ Automat.\ Control}, vol.\ 57, no.\ 4, pp.\ 1035--1040, 2012.

\bibitem{rawlings2017}
J.\ B.\ Rawlings, D.\ Q.\ Mayne, and M.\ Diehl,
\emph{Model Predictive Control: Theory, Computation, and Design}, 2nd~ed.,
Nob~Hill~Publishing, 2017.

\bibitem{utkin1992}
V.\ I.\ Utkin,
\emph{Sliding Modes in Control and Optimization},
Springer-Verlag, 1992.

\bibitem{slotine1991}
J.-J.\ E.\ Slotine and W.\ Li,
\emph{Applied Nonlinear Control},
Prentice Hall, 1991.

\end{thebibliography}

\end{document}
